% Reciprocal proofs for ED26-REQ-N-c-002-02, Fagin 1976.

\section{Fagin's finite-model probabilities}
\label{sec:annex-p09-n-c-002-02}

\begin{annexitem}{annex:p9.n-c-002.2.extension-estimate}{thm:p09-fagin1976-extension-estimate}{Counting fresh witnesses for every finite positive arity}
\begin{proof}
Fix a distinct ordered parameter tuple \(\bar a=(a_1,\ldots,a_m)\)
and condition on its induced open description \(M\).  For each
\(i\)-ary symbol, there are \((m+1)^i\) ordered tuples on the
parameter points and one new point, of which \(m^i\) use only
parameters.  The difference counts exactly the atoms involving the new
point, including atoms with several occurrences of that point.  Because
the points are distinct, different variable tuples give different ground
relation facts.  The new facts are independent of all parameter-only
facts, and of each other.  A specified extending description therefore
has probability \(q=2^{-d_Y(m)}\) at a given candidate \(b\notin\bar a\).

For two distinct candidates \(b,c\notin\bar a\), the atoms inspected
for \(b\) have entries only from \(\bar a\cup\{b\}\), contain
\(b\), and do not contain \(c\).  Those inspected for \(c\) contain
\(c\) and not \(b\).  The two sets of atoms are disjoint even at
arity greater than two.  No atom involving both candidates is inspected.
This proves independence of all \(n-m\) candidate events, including
after conditioning on \(M\).  The conditional probability of no witness
is consequently \((1-q)^{n-m}\).  Multiplying by the probability of
\(M\) can only reduce this bound.  A tuple not satisfying \(M\) is
not a counterinstance to the axiom.

There are \(n(n-1)\cdots(n-m+1)\leq n^m\) ordered distinct
parameter tuples.  A union bound over them gives the claimed estimate;
independence between different parameter tuples is neither asserted nor
needed.  If \(0<q<1\), put \(c=-\log(1-q)>0\).  The logarithm
of the upper bound is \(m\log n-c(n-m)\), which tends to minus
infinity because \(\log n/n\to0\).  If \(q=1\), no tuple
can lack a witness when \(n>m\), so the probability is already zero.
This includes the empty relation vocabulary.  The argument expands the
source's displayed binary estimate to arbitrary finite arities
\cite[thm.~2, p.~52]{Fagin1976}.
\end{proof}
\end{annexitem}

\begin{annexitem}{annex:p9.n-c-002.2.theory}{thm:p09-fagin1976-theory}{Existence, both back-and-forth directions, and the completeness test}
\begin{proof}
For a finite \(T_0=\{e_1,\ldots,e_t\}\subseteq T\), the extension
estimate gives
\(p_n(\bigwedge_j e_j)\geq1-\sum_jp_n(\neg e_j)>0\)
for all sufficiently large \(n\).  Thus every finite subset has a
finite model.  First-order compactness gives a model of \(T\).
No model of \(T\) is finite: if a model had exactly \(m\) elements,
list all of them as distinct parameters, take their actual description
\(M\), and choose any extending description \(N\).  The corresponding
axiom requires an element outside the entire domain.  This is impossible.
Downward L\"owenheim--Skolem in the finite vocabulary supplies a countably
infinite elementary substructure of a model, hence a countable model of
\(T\).  These uses of compactness and downward L\"owenheim--Skolem
are prerequisites, not finite-model compactness claims.

Let \(\mathcal A\) and \(\mathcal B\) be countably infinite models of \(T\), with
enumerations \(a_0,a_1,\ldots\) and \(b_0,b_1,\ldots\).
Start with the empty partial bijection.  Suppose a finite partial
isomorphism \(g\) maps distinct \(u_1,\ldots,u_m\) to
\(v_1,\ldots,v_m\).  At an even stage choose the first \(a_j\)
not in its domain.  Let \(M\) be the actual description of
\(\bar u\) in \(\mathcal A\), and \(N\) that of
\((\bar u,a_j)\).  Since \(g\) is a partial isomorphism,
\(\bar v\) satisfies \(M\) in \(\mathcal B\).  Its extension
axiom gives a fresh \(b\) with description \(N(\bar v,b)\).
The map \(g\cup\{(a_j,b)\}\) is a bijection between the enlarged
finite sets, preserves equality by freshness, and preserves all relation
facts by the description.  In particular repeated arguments and high
arities are preserved, not just binary edges.  At an odd stage choose the
first \(b_j\) outside the range and apply exactly this argument to
\(g^{-1}\); then invert the extended map.

The union of the finite maps is injective since every stage is a
bijection on its finite sets, total by the even stages, and onto by the
odd stages.  Every relation tuple is finite and eventually lies in the
domain of one stage, where its truth and falsity are preserved.  The
union is therefore an isomorphism.  This proves countable categoricity
without a name or computable coding for the distinguished model.

For completeness, suppose neither \(T\entails\sigma\) nor
\(T\entails\neg\sigma\).  By the definition of semantic consequence,
there are models of \(T\cup\{\neg\sigma\}\) and of
\(T\cup\{\sigma\}\).  Both are infinite.  Downward
L\"owenheim--Skolem gives countably infinite elementary substructures,
still satisfying their respective sentences.  Countable categoricity
makes these substructures isomorphic; isomorphism invariance of sentence
truth contradicts their opposite truth values for \(\sigma\).
Consistency excludes both consequences holding together.  This is the
\L o\'s--Vaught argument expanded from Fagin's explicitly sketched
Theorem~1 \cite[p.~51]{Fagin1976}.
\end{proof}
\end{annexitem}

\begin{annexitem}{annex:p9.n-c-002.2.labelled}{thm:p09-fagin1976-labelled}{The finite logical bridge and geometric convergence}
\begin{proof}
If \(T\entails\sigma\), compactness gives a finite
\(\{e_1,\ldots,e_t\}\subseteq T\) entailing \(\sigma\).
Every countermodel of \(\sigma\) must violate at least one \(e_j\).
Hence \(p_n(\neg\sigma)\leq\sum_jp_n(\neg e_j)\to0\), so
\(p_n(\sigma)\to1\).  The conclusion also holds when the finite
set is empty, in which case \(\sigma\) is valid.  Completeness
gives exactly one of \(T\entails\sigma,T\entails\neg\sigma\).
The second case gives \(p_n(\neg\sigma)\to1\) and
\(p_n(\sigma)\to0\).  Conversely, if \(p_n(\sigma)\to1\)
the second case is impossible.  This proves both the zero--one law
and the asserted equivalence.

For the rate, a nonzero individual error bound has form
\(n^m k^{n-m}=k^{-m}n^m k^n\), with \(0<k<1\).
Choose \(k<\theta_0<1\).  The ratio of this bound to
\(\theta_0^n\) is \(k^{-m}n^m(k/\theta_0)^n\to0\)
by the same logarithmic estimate as before.  Therefore each individual
error is eventually less than \(\theta_0^n\).  For the finitely
many axioms, take the largest chosen \(\theta_0\), still less than
one, and then choose \(\theta_0<\theta<1\).  The finite sum is
at most \(t\theta_0^n<\theta^n\) eventually.  An identically
zero error needs no bound.  Complementation handles limit zero and
\(\lvert1-p_n(\sigma)\rvert\) handles limit one.  The sentence
is fixed before these constants are chosen; no uniform rate over all
sentences is proved \cite[cor.~3, thm.~4, and thm.~10]{Fagin1976}.
\end{proof}
\end{annexitem}

\begin{annexitem}{annex:p9.n-c-002.2.finite-fragments}{thm:p09-fagin1976-finite-fragments}{Finite fragments and the original-vocabulary spectrum argument}
\begin{proof}
For a nonempty finite subset \(\{e_1,\ldots,e_t\}\) choose \(N\)
so that every \(p_n(\neg e_j)<1/t\) when \(n>N\).
Then \(p_n(\bigwedge_j e_j)>0\).  A positive proportion means an
actual structure of cardinality \(n\) exists.  The empty subset is
satisfied by every structure, so it needs no division by \(t\).

For the spectrum claim in a finite vocabulary with functions and
constants, take the following theory \(\Sigma\).  It asserts at
least \(k\) distinct elements for every positive integer \(k\);
each \(r\)-ary relation is full; each positive-arity function satisfies
\(\forall x_1\cdots x_r\,F(x_1,\ldots,x_r)=x_1\); and every
two constants are equal.  Nullary function symbols are constants here.
If nullary predicates are permitted, assign them true as well.
The theory has an infinite model: on any infinite set choose an element
for all constants, use the prescribed projections, and make relations
full.  Any two countably infinite such models are isomorphic: map the
common constant value to the other one when constants exist and extend
to a bijection of the remaining countable sets.  Every bijection so
chosen preserves the full relations and projection functions.  If there
are no constants, any bijection works.  There are no finite models.
The same downward L\"owenheim--Skolem argument used in the preceding
annex item therefore proves \(\Sigma\) complete.

Exactly one of \(\Sigma\entails\sigma\) and
\(\Sigma\entails\neg\sigma\) follows.  Denote it by
\(\Sigma\entails\rho\).  Compactness gives a finite
\(\Sigma_0\subseteq\Sigma\) entailing \(\rho\).
Choose \(N\) greater than every lower cardinality bound occurring in
\(\Sigma_0\) (or choose \(N=1\) if none occurs).  For every
\(n>N\), the full-relation/projection structure on \([n]\), with
all constants interpreted by \(1\), satisfies \(\Sigma_0\).
It thus satisfies \(\rho\).  This proves the cofinite spectrum of
\(\rho\) in the original vocabulary, as in the second proof of
Fagin's Corollary~6.  No uniform random-function argument was used.

To locate the first proof's boundary precisely, a standard relational
encoding of a function sentence conjoins a condition that its graph
relation is total and single-valued.  Negating that encoded sentence
allows the coding condition to fail, independently of the truth of the
original negated sentence.  Thus equality of the encoded and original
\emph{positive} spectra by itself is insufficient for the negation
step.  The canonical-structure proof above avoids that inference and
proves the actual stated disjunction \cite[pp.~53--54]{Fagin1976}.
\end{proof}
\end{annexitem}

\begin{annexitem}{annex:p9.n-c-002.2.monadic}{thm:p09-fagin1976-monadic}{Stars and bars, the threshold induction, and monadic probabilities}
\begin{proof}
List all \(s=2^h\) one-point types.  Unary relations and equality
have no facts relating distinct points beyond their type memberships,
so a bijection preserving the type of each point is an isomorphism.
It exists exactly when the type counts agree.  To count nonnegative
\(s\)-tuples summing to \(n\), place \(n\) stars and \(s-1\)
bars in a line.  The numbers of stars before, between, and after the
bars give the tuple, including zero entries and consecutive bars.
Conversely each tuple gives one line.  Selecting the positions of the
bars gives \(\binom{n+s-1}{s-1}\).  If one count is fixed at
\(k<n\) and \(s\geq2\), repeat this argument on the remaining
\(s-1\) counts with sum \(n-k\), yielding the stated numerator.
Its degree in \(n\) is \(s-2\), whereas the denominator has
degree \(s-1\), so the ratio tends to zero.

Here is the logical threshold argument, including formulas with free
variables.  Let \(g\) match finitely many named points bijectively
between two monadic structures, preserving their types and all
equalities of assigned variables.  Suppose that, in each type, the
numbers of \emph{unmatched} points on the two sides are equal after
truncation at \(t\).  Every formula of quantifier rank at most
\(t\), with its free variables assigned to the matched points,
has the same truth value on both sides.  Prove this by structural
induction on the formula, retaining the indicated rank bound.
Atomic formulas agree by the partial bijection and point types;
Boolean operations preserve agreement.  At an existential quantifier
with \(t\geq1\), a witness already matched has its image as witness.
For a fresh witness of a given type, that type has an unmatched point
on the other side: if its count is below \(t\), the counts are
equal; otherwise both are at least \(t\).  Match the witness to
one such point.  After removing one point on each side in this type,
the remaining counts agree after truncation at \(t-1\): equal
counts below \(t\) decrease equally, and two counts at least
\(t\) remain at least \(t-1\).  Other types also retain agreement
at the smaller threshold.  The induction applies to the matrix.
The converse existential direction uses \(g^{-1}\).
Universal quantifiers follow by negation.  This proves the claim.

For a rank-\(r\) sentence start with the empty matching.  It depends
only on the finite vector \((\min(c_i,r))_{i=1}^s\).
Each vector entry \(j<r\) is specified by ``exactly \(j\) points
of type \(i\)'', and entry \(r\) by negating the finitely many
``exactly \(j\)'' assertions with \(j<r\).
Take the disjunction of the vector descriptions on which the sentence
is true.  Unattainable vectors may be omitted; if no vector is true,
use the empty disjunction, namely false.  This yields the claimed
Boolean combination.  Rank zero sentences are constants, since the
vocabulary has no constants or nullary predicates.

For \(s\geq2\), each event of exactly \(j<r\) points of a type
has unlabelled probability tending to zero by the counted ratio.
Under labelled sampling it has probability
\[
 \binom nj(1/s)^j(1-1/s)^{n-j}\longrightarrow0,
\]
because each of the \(n\) points independently has each type with
probability \(1/s\).  A union bound makes the probability of any
of these finitely many exceptional events tend to zero in either
measure.  Outside their union every type has at least \(r\) members;
the sentence has one fixed truth value there by the threshold claim.
Hence the two limits exist, are zero or one, and agree.

For the rational-function refinement, partition the count vectors into
their threshold descriptions.  Suppose a description fixes \(j\)
counts below \(r\) and leaves \(t=s-j\) counts at least \(r\).
Subtract \(r\) from each latter count.  Their remaining sum is
\(n-K\) for a fixed integer \(K\).  For \(t\geq1\), the
number of such vectors is \(\binom{n-K+t-1}{t-1}\) for large
\(n\), a polynomial in \(n\).  For \(t=0\) the description
is impossible at all sufficiently large \(n\).  Sum these
polynomials over the finitely many descriptions satisfying the sentence
and divide by \(b_n(Y)\).  This proves eventual rationality.
If \(s=1\), the only count is \(n\); its truncated value is
\(r\) for all \(n\geq r\), so equality-only sentences are
eventually constant.  The source's monadic argument, including its
external stars-and-bars reference, is thus supplied explicitly
\cite[lem.~7 and thm.~9, case~1]{Fagin1976}.
\end{proof}
\end{annexitem}

\begin{annexitem}{annex:p9.n-c-002.2.unlabelled}{thm:p09-fagin1976-unlabelled}{The nonmonadic transfer and its precise imported premise}
\begin{proof}
The monadic case was proved in the preceding lemma.  Suppose now that
\(Y\) is nonmonadic and \(p_n(\sigma)\to1\).
If there are \(x_n\) labelled models and \(y_n\) model types of
size \(n\), each type accounts for at most \(n!\) labellings,
so \(y_n\geq x_n/n!\).  Since
\(x_n=p_n(\sigma)a_n(Y)\) and
\(y_n=v_n^Y(\sigma)b_n(Y)\),
\[
 v_n^Y(\sigma)\geq
 p_n(\sigma)\frac{a_n(Y)}{b_n(Y)n!}.
\]
The right side tends to one by the cited enumeration estimate, and
the left side is at most one, so \(v_n^Y(\sigma)\to1\).
If \(p_n(\sigma)\to0\), apply the same argument to
\(\neg\sigma\) and complement.  Notice that the bound does
not replace either measure by the other at a finite cardinality.

For the rate set \(r_n=a_n(Y)/(b_n(Y)n!)\).
Always \(0<r_n\leq1\), since \(a_n(Y)\leq b_n(Y)n!\).
In the limit-one case the same inequality gives
\[
 0\leq1-v_n^Y(\sigma)
 \leq 1-p_n(\sigma)r_n
 =1-p_n(\sigma)+p_n(\sigma)(1-r_n)
 \leq (1-p_n(\sigma))+(1-r_n).
\]
Both summands converge geometrically, the first by the labelled law
and the second by the expressly imported rate.  A finite-sum estimate
as in that law gives a geometric bound for their sum.  Complementation
again handles limit zero.  This proves the source's Theorems~9--10
\emph{using}, not reproving, Theorem~8 and the rate cited from its
reference~[3] \cite[pp.~54--56]{Fagin1976}.

With one unary \(U\), the possible counts \(0,\ldots,n\) are
the \(n+1\) types of structures; only count \(n\) satisfies
\(\forall x\,U(x)\), giving \(1/(n+1)\).  This is not
geometrically fast: for every \(0<\theta<1\),
\((n+1)\theta^n\to0\), so eventually \(1/(n+1)>\theta^n\).
Labelled probability is \(2^{-n}\) by the \(n\) independent
unary facts.  For binary \(P\), the diagonal assertion likewise
has labelled probability \(2^{-n}\); the already proved
nonmonadic transfer gives its geometric unlabelled rate.
\end{proof}
\end{annexitem}

\begin{annexitem}{annex:p9.n-c-002.2.graph}{thm:p09-fagin1976-graph}{Graph and tournament extensions with fixed independent edge choices}
\begin{proof}
Fix disjoint parameter sets \(U,V\) of total size \(m=u+v\).
A candidate \(z\) outside them inspects the \(m\) unordered
pairs joining \(z\) to the parameters.  Their choices are independent,
so success has probability \(p^u(1-p)^v\).  Distinct outside
candidates inspect disjoint sets of unordered pairs: they never inspect
the pair joining them.  Their success events are therefore independent.
The exact no-witness probability for the fixed parameters is
\((1-q)^{n-m}\).  There are at most \(n^m\) ordered choices
of the parameters, giving the union bound.  Since \(p\) is fixed
strictly between zero and one, \(q>0\) is fixed, and the bound
tends to zero by the previous logarithmic calculation.  The case
\(m=0\) holds automatically in nonempty graphs.

Every finite set of graph extension axioms thus has a finite graph model
by a finite union bound.  Compactness gives a model of the full theory,
and the extension requirements rule out finite models by requesting a
fresh point outside all their vertices.  Downward L\"owenheim--Skolem
gives a countable model.  To extend a finite partial graph isomorphism,
partition its domain into neighbors and nonneighbors of the desired new
vertex.  Their images give \(U,V\) on the other side; the graph
extension axiom gives a fresh image with precisely the needed incidences.
Symmetry fixes the reverse incidences, and irreflexivity fixes the loop.
Alternating domain and range extensions along enumerations gives an
isomorphism exactly as in the general theory proof.  Completeness follows
by the same elementary-substructure argument.  For any fixed sentence,
compactness yields finitely many extension axioms entailing it or its
negation, and the finite union bound yields the indicated limit.
All fixed \(p\in(0,1)\) have the same extension theory, hence the
same limit for each sentence.  This extends the source's conditioned
fair-graph calculation without asserting a varying-density theorem.

For \(u=1,v=0\), every vertex has a fresh neighbor outside itself,
so the complete graph extension theory entails
\(\forall x\exists y E(x,y)\).  Its negation is the isolated-vertex
sentence in the body, giving the two stated limits.  The numerical
\(u=v=1,p=1/2\) example substitutes \(q=1/4\) in the bound.

Under the original unrestricted binary distribution, conditioning on
\(\tau_1\) leaves exactly \(2^{\binom n2}\) structures, one for
each choice of unordered edges.  They originally had equal weight, so
the conditional distribution is uniform on these choices.  Its
unordered-edge variables are independent and fair.  Conditioning instead
on \(\tau_2\) leaves exactly \(2^{\binom n2}\) tournaments, with
one of the two orientations chosen independently and fairly for each
unordered pair.  For any fixed oriented one-point extension on \(m\)
parameters, a candidate has probability \(2^{-m}\) of the required
directions.  Distinct candidates again inspect disjoint pairs, yielding
the bound \(n^m(1-2^{-m})^{n-m}\).  For an extension of a finite
partial tournament isomorphism, list the required direction between the
new domain point and each old domain point.  The corresponding one-point
axiom in the range gives a fresh point with those directions to the
images.  The opposite directions are forced by \(\tau_2\), and
the false loop is also forced.  Thus the enlarged map preserves all
binary facts.  The reverse step extends the inverse map in the same way.
Alternation along the enumerations consequently gives a total onto
isomorphism by the already proved union argument.  A finite union of
extension failures has probability tending to zero in the fair
tournament space, so finite fragments have finite models.  Compactness
gives existence of the full extension theory; the fresh-point
requirement gives infinitude, and downward L\"owenheim--Skolem and
the proved countable uniqueness give completeness.  A final compactness
step and finite union bound give the conditional sentence law.  This
supplies both conditional zero--one examples in section~4(6)
\cite[p.~57]{Fagin1976}.
\end{proof}
\end{annexitem}

\begin{annexitem}{annex:p9.n-c-002.2.boundaries}{thm:p09-fagin1976-boundaries}{Pairings, constant and function counts, and the conditional failures}
\begin{proof}
In the second-order example define, from a proposed relation \(P\),
the relation \(H(x,y)\) by
\(x\ne y\land P(x,y)\land P(y,x)\).  It is symmetric and
irreflexive, and the matrix asserts that each point has exactly one
\(H\)-neighbor.  The neighbor operation \(g\) therefore satisfies
\(g(x)\ne x\) and \(g(g(x))=x\): symmetry makes \(x\) a
neighbor of \(g(x)\), and uniqueness forces it to be that point's
sole neighbor.  The finite universe splits into disjoint two-element
pairs, so its cardinality is even.  Conversely, pair the points of any
even universe and let \(P\) hold in both directions exactly on each
pair.  The matrix is then true.  At every positive cardinality the
empty vocabulary has a single labelled structure and a single
isomorphism type, hence the two truth probabilities are precisely the
even-cardinality indicator, with no limit.  The quantified relation
is a second-order witness, not a sampled extra symbol.

For \(U,c\), there are \(n2^n\) labelled structures: independently
choose a subset for \(U\) and one of \(n\) values for \(c\).
For each fixed value, exactly \(2^{n-1}\) subsets contain it, so
the favourable count is \(n2^{n-1}\) and the fraction is \(1/2\).
For a unary \(F\), the \(n\) function values are independent
uniform choices from \([n]\).  Avoiding the fixed point at each
argument permits \(n-1\) values, giving \((n-1)^n\) favourable
maps among \(n^n\) maps.  Since
\(n\log(1-1/n)\to-1\), the fraction tends to \(e^{-1}\).
At \(n=1\) it is zero, which does not affect the limit.

The same pairing argument for the first-order \(\sigma_1\)
shows that it has a model exactly at even sizes.  Here \(P\) is
part of a sampled structure rather than a quantified witness, but all
structures satisfying \(\sigma_1\) still induce such a pairing.
For \(\sigma_2\) choose its witness \(w\); the reciprocal
\(P\)-edges on the other \(n-1\) points form a pairing, so
\(n\) is odd.  Conversely pair those \(n-1\) points and
leave \(w\) without reciprocal edges.  This gives a model at
every odd size, including \(n=1\), where the universal condition
outside \(w\) is vacuous.  Nonreciprocal extra \(P\)-edges are
irrelevant and do not change either parity argument.

Thus \(\tau=\sigma_1\lor\sigma_2\) has a model at every
size and therefore has positive probability in each of the two finite
uniform spaces.  At even sizes every model of \(\tau\) satisfies
\(\sigma_1\), while at odd sizes none does.  Division by the
positive probability of \(\tau\) gives exactly the alternating
conditional probabilities in the statement.  This shows why the
unconditional law cannot be divided by an arbitrary event whose
probability may tend to zero.

Finally \(\tau_3\) leaves just two unary structures on \([n]\):
\(U\) is full or empty.  Both have equal conditional labelled
weight.  In either structure a first-order sentence reduces, by
replacing its unary atoms by true or false, to a sentence using only
equality.  The monadic threshold argument with one point type shows
that the truth of each resulting sentence is constant for all large
\(n\).  The conditional probability is the mean of these two
eventually constant indicators, and so is eventually \(0,1/2\),
or \(1\).  The full-\(U\) sentence has one true indicator and
one false indicator at every nonempty size, proving its \(1/2\)
value \cite[sec.~4(1), (2), (5), and (6)]{Fagin1976}.
\end{proof}
\end{annexitem}
